\documentclass[10pt,a4paper]{amsart}
\usepackage{amsmath,amssymb,amsthm,graphicx}
\usepackage{verbatim}
\newenvironment{reference}{\comment}{\endcomment}
%\newenvironment{reference}{}{}
\newenvironment{slogan}{\comment}{\endcomment}
\newenvironment{history}{\comment}{\endcomment}
% For commutative diagrams we use Xy-pic
\usepackage[all]{xy}

% We use 2cell for 2-commutative diagrams.
\xyoption{2cell}
% Theorem environments.
%
\theoremstyle{plain}
\newtheorem{theorem}[subsection]{Theorem}
\newtheorem{proposition}[subsection]{Proposition}
\newtheorem{lemma}[subsection]{Lemma}

\theoremstyle{definition}
\newtheorem{definition}[subsection]{Definition}
\newtheorem{example}[subsection]{Example}
\newtheorem{exercise}[subsection]{Exercise}
\newtheorem{situation}[subsection]{Situation}

\theoremstyle{remark}
\newtheorem{remark}[subsection]{Remark}
\newtheorem{remarks}[subsection]{Remarks}

\numberwithin{equation}{subsection}

% Macros
%
\def\lim{\mathop{\mathrm{lim}}\nolimits}
\def\colim{\mathop{\mathrm{colim}}\nolimits}
\def\Spec{\mathop{\mathrm{Spec}}}
\def\Hom{\mathop{\mathrm{Hom}}\nolimits}
\def\Ext{\mathop{\mathrm{Ext}}\nolimits}
\def\SheafHom{\mathop{\mathcal{H}\!\mathit{om}}\nolimits}
\def\SheafExt{\mathop{\mathcal{E}\!\mathit{xt}}\nolimits}
\def\Sch{\mathit{Sch}}
\def\Mor{\mathop{\mathrm{Mor}}\nolimits}
\def\Ob{\mathop{\mathrm{Ob}}\nolimits}
\def\Sh{\mathop{\mathit{Sh}}\nolimits}
\def\NL{\mathop{N\!L}\nolimits}
\def\CH{\mathop{\mathrm{CH}}\nolimits}
\def\proetale{{pro\text{-}\acute{e}tale}}
\def\etale{{\acute{e}tale}}
\def\QCoh{\mathit{QCoh}}
\def\Ker{\mathop{\mathrm{Ker}}}
\def\Im{\mathop{\mathrm{Im}}}
\def\Coker{\mathop{\mathrm{Coker}}}
\def\Coim{\mathop{\mathrm{Coim}}}

% Boxtimes
%
\DeclareMathSymbol{\boxtimes}{\mathbin}{AMSa}{"02}

%
% Macros for moduli stacks/spaces
%
\def\QCohstack{\mathcal{QC}\!\mathit{oh}}
\def\Cohstack{\mathcal{C}\!\mathit{oh}}
\def\Spacesstack{\mathcal{S}\!\mathit{paces}}
\def\Quotfunctor{\mathrm{Quot}}
\def\Hilbfunctor{\mathrm{Hilb}}
\def\Curvesstack{\mathcal{C}\!\mathit{urves}}
\def\Polarizedstack{\mathcal{P}\!\mathit{olarized}}
\def\Complexesstack{\mathcal{C}\!\mathit{omplexes}}
% \Pic is the operator that assigns to X its picard group, usage \Pic(X)
% \Picardstack_{X/B} denotes the Picard stack of X over B
% \Picardfunctor_{X/B} denotes the Picard functor of X over B
\def\Pic{\mathop{\mathrm{Pic}}\nolimits}
\def\Picardstack{\mathcal{P}\!\mathit{ic}}
\def\Picardfunctor{\mathrm{Pic}}
\def\Deformationcategory{\mathcal{D}\!\mathit{ef}}
\usepackage{cleveref}
\crefname{originalreference}{original reference}{original references}
\Crefname{originalreference}{Original reference}{Original references}
\makeatletter
\newlabel{criteria-lemma-flat-quotient-flat-presentation}{{criteria-lemma-flat-quotient-flat-presentation}{source}{}{sourcepointer-1}{}}
\newlabel{criteria-lemma-flat-quotient-flat-presentation@cref}{{[originalreference][0][]criteria-lemma-flat-quotient-flat-presentation}{[1][0][]source}{}{}{}}
\makeatother
\crefname{originalreference}{original reference}{original references}
\Crefname{originalreference}{Original reference}{Original references}
\makeatletter
\newlabel{spaces-groupoids-lemma-quotient-stack-2-cartesian}{{spaces-groupoids-lemma-quotient-stack-2-cartesian}{source}{}{sourcepointer-2}{}}
\newlabel{spaces-groupoids-lemma-quotient-stack-2-cartesian@cref}{{[originalreference][0][]spaces-groupoids-lemma-quotient-stack-2-cartesian}{[1][0][]source}{}{}{}}
\makeatother
\crefname{originalreference}{original reference}{original references}
\Crefname{originalreference}{Original reference}{Original references}
\makeatletter
\newlabel{stacks-properties-lemma-characterize-surjective}{{stacks-properties-lemma-characterize-surjective}{source}{}{sourcepointer-3}{}}
\newlabel{stacks-properties-lemma-characterize-surjective@cref}{{[originalreference][0][]stacks-properties-lemma-characterize-surjective}{[1][0][]source}{}{}{}}
\makeatother
\crefname{originalreference}{original reference}{original references}
\Crefname{originalreference}{Original reference}{Original references}
\makeatletter
\newlabel{stacks-properties-lemma-points-cartesian}{{stacks-properties-lemma-points-cartesian}{source}{}{sourcepointer-4}{}}
\newlabel{stacks-properties-lemma-points-cartesian@cref}{{[originalreference][0][]stacks-properties-lemma-points-cartesian}{[1][0][]source}{}{}{}}
\makeatother
\title{The space of a quotient stack is the quotient of the space of the covering by the relation coming from the groupoid}
\author{The Stacks Project Authors}
\date{stacks-project:stacks-morphisms.tex:local-snapshot-20261012; GNU FDL 1.2 or later}
\begin{document}
\maketitle
\noindent\textit{Editorial note.} All delivered mathematical text is the author's own, copied from the pinned Stacks Project chapter source sources/stacks-algebraic-stacks-morphisms/pinned-original-source.tex; no statement or proof step was rewritten, reordered or removed.\par\medskip
\noindent\textit{Editorial note.} Preamble provenance: the Stacks Project keeps one single shared preamble for every chapter; this card takes it directly from the staged sources/stacks-algebraic-stacks-morphisms/preamble.tex as three verbatim spans (lines 8-12, 14-18 and 154-217), so no derived file was written and nothing was added or redefined. The wrapper class and the base packages are supplied by the assembler.\par\medskip
\noindent\textit{Editorial note.} Bibliography provenance: neither the counted statement nor its proof contains a citation (this lane rescans the whole span before writing the card), so the card carries no bibliography and no entry was reconstructed.\par\medskip
\noindent\textit{Editorial note.} Selection provenance: this card was opened under the file-based control contract (CONTROL.md), with the mechanical pre-card gate run on its statement and proof interval before the card was written.\par\medskip
\noindent\textit{Allowed standard prerequisites (editorial).}\begin{itemize}
\item The statement that the quotient of a flat groupoid in algebraic spaces is an algebraic stack with surjective flat and locally of finite presentation structure morphism. Invoked inside the counted proof (source lines 4135-4138) as a cross-chapter pointer. Makes the quotient morphism from the covering space onto the stack.
\item The characterization of surjectivity of a morphism of algebraic stacks by surjectivity on points. Invoked inside the counted proof (source lines 4139-4141) as a cross-chapter pointer. Turns surjectivity of the quotient morphism into surjectivity of the map on spaces.
\item The identification of the groupoid space with the fibre product of the quotient morphism with itself. Invoked inside the counted proof (source lines 4142-4144) as a cross-chapter pointer. Describes the relation as a fibre product.
\item The statement describing the points of a fibre product of algebraic stacks. Invoked at the end of the counted proof (source lines 4145-4149) as a cross-chapter pointer. Identifies the image of the groupoid space in the square of the covering with the pairs having the same image in the stack.
\end{itemize}\medskip
\ifdefined\hypertarget\hypertarget{sourcepointer-1}{}\fi
\noindent\textit{Source reference.} Original reference criteria-lemma-flat-quotient-flat-presentation: sources/stacks-algebraic-stacks-morphisms/pinned-original-source.tex, source lines 3877--3877. Navigation pointer only: the counted text refers to the Stacks item criteria-lemma-flat-quotient-flat-presentation, which belongs to another chapter of the project and is not part of this card. The pointer carries the target's own label as its number because this chapter's staged snapshot does not contain the other chapter's tag table; its content is not reproduced here.\par\medskip
\ifdefined\hypertarget\hypertarget{sourcepointer-2}{}\fi
\noindent\textit{Source reference.} Original reference spaces-groupoids-lemma-quotient-stack-2-cartesian: sources/stacks-algebraic-stacks-morphisms/pinned-original-source.tex, source lines 4107--4107. Navigation pointer only: the counted text refers to the Stacks item spaces-groupoids-lemma-quotient-stack-2-cartesian, which belongs to another chapter of the project and is not part of this card. The pointer carries the target's own label as its number because this chapter's staged snapshot does not contain the other chapter's tag table; its content is not reproduced here.\par\medskip
\ifdefined\hypertarget\hypertarget{sourcepointer-3}{}\fi
\noindent\textit{Source reference.} Original reference stacks-properties-lemma-characterize-surjective: sources/stacks-algebraic-stacks-morphisms/pinned-original-source.tex, source lines 4141--4141. Navigation pointer only: the counted text refers to the Stacks item stacks-properties-lemma-characterize-surjective, which belongs to another chapter of the project and is not part of this card. The pointer carries the target's own label as its number because this chapter's staged snapshot does not contain the other chapter's tag table; its content is not reproduced here.\par\medskip
\ifdefined\hypertarget\hypertarget{sourcepointer-4}{}\fi
\noindent\textit{Source reference.} Original reference stacks-properties-lemma-points-cartesian: sources/stacks-algebraic-stacks-morphisms/pinned-original-source.tex, source lines 2126--2126. Navigation pointer only: the counted text refers to the Stacks item stacks-properties-lemma-points-cartesian, which belongs to another chapter of the project and is not part of this card. The pointer carries the target's own label as its number because this chapter's staged snapshot does not contain the other chapter's tag table; its content is not reproduced here.\par\medskip
\begin{lemma}
\label{lemma-points-presentation}
Let $(U, R, s, t, c)$ be a groupoid in algebraic spaces such that
$s, t : R \to U$ are flat and locally of finite presentation.
Consider the algebraic stack $\mathcal{X} = [U/R]$ (see above).
Then the image of $|R| \to |U| \times |U|$ is an equivalence relation
and $|\mathcal{X}|$ is the quotient of $|U|$ by this equivalence relation.
\end{lemma}

\begin{proof}
The induced morphism $p : U \to \mathcal{X}$ is surjective, flat,
and locally of finite presentation, see
Criteria for Representability, Lemma
\ref{criteria-lemma-flat-quotient-flat-presentation}.
Hence $|U| \to |\mathcal{X}|$ is surjective by
Properties of Stacks, Lemma
\ref{stacks-properties-lemma-characterize-surjective}.
Note that $R = U \times_\mathcal{X} U$, see
Groupoids in Spaces,
Lemma \ref{spaces-groupoids-lemma-quotient-stack-2-cartesian}.
Hence
Properties of Stacks, Lemma \ref{stacks-properties-lemma-points-cartesian}
implies the map
$$
|R| \longrightarrow |U| \times_{|\mathcal{X}|} |U|
$$
is surjective. Hence the image of $|R| \to |U| \times |U|$ is
exactly the set of pairs $(u_1, u_2) \in |U| \times |U|$
such that $u_1$ and $u_2$ have the same image in $|\mathcal{X}|$.
Combining these two statements we get the result of the lemma.
\end{proof}

\end{document}
